\section{Directed relations without transitive triangles}\label{sec:directed}

By \cref{lem:tt-paired}, a loopless relation is TT-free exactly when every coordinate has paired distance at least
five. This section proves the half bound for three classes of such relations: subcubic underlying graphs
(Theorem~U3), relations whose vertices of in-degree at least three are acyclic (Theorem~T), and block structures with
private outside parents (Theorem~PC), which contain all TT-free Cayley relations $\Cay(G,\{s,s^{-1},t\})$
(Corollary~CC). Directed triangles are allowed throughout; they are the only triangles a TT-free relation has.

\begin{lemma}[Triangles of TT-free relations]\label{lem:dtri}
\formal{U3.lean}{dtri_orient}
Let $N$ be TT-free, and let $x_0,x_1,x_2$ be distinct and pairwise adjacent in the underlying graph, with
$x_0\to x_1$. Then $x_1\to x_2$ and $x_2\to x_0$, and none of the arcs $x_1\to x_0$, $x_2\to x_1$, $x_0\to x_2$ is
present.
\end{lemma}

\begin{outline}
If $x_2\to x_1$, then either $x_0\to x_2$, and $x_0\to x_2\to x_1$ with $x_0\to x_1$ is a transitive triangle, or
$x_2\to x_0$, and $x_2\to x_0\to x_1$ with $x_2\to x_1$ is one. So $x_1\to x_2$, and then $x_0\to x_2$ would close
the transitive triangle $x_0\to x_1\to x_2$; hence $x_2\to x_0$. The reverse arcs $x_1\to x_0$ and $x_2\to x_1$ would
close the transitive triangles $x_1\to x_2\to x_0$, $x_1\to x_0$ and $x_2\to x_0\to x_1$, $x_2\to x_1$.
\end{outline}

\subsection{Subcubic underlying graphs: Theorem U3}\label{sec:U3}

\begin{theorem}[Theorem U3]\label{thm:U3}
\formal{U3.lean}{und, sub_nbhd_und, dtriPiece, triPiece, rateBound_u3, card_sq_le_u3}
Let $N$ be loopless and TT-free, and let its underlying graph have maximum degree at most~$3$. Then $N$ satisfies
\eqref{eq:star}, and every storage code $C$ for $N$ over a nonempty finite alphabet satisfies
$\card C^2\le\card\alpha^{\card V}$.
\end{theorem}

\begin{outline}
Let $G=U(N)$; since $N$ is loopless, $N(v)\subseteq N_G(v)$ (\lean{sub_nbhd_und}). Run the shortest-cycle decomposition
of Theorem~N3 on $G$ (\lean{decomp_of_cycles}). Shortest cycles of length at least four are pieces for $N_G$ and hence
for $N$, which has fewer arcs (\cref{lem:greedy}\ref{it:mono}); the acyclic remainder is a relative $2$-colouring, built
directly for $N$ by \cref{lem:Bdir}. The new case is a shortest
cycle of length three. By \cref{lem:dtri} it is a directed $3$-cycle $c_0\to c_1\to c_2\to c_0$ without reverse arcs
(\lean{triPiece}). Take $A=\{c_0,c_2\}$, which carries the single arc $c_2\to c_0$ and is acyclic, and $B=\{c_1\}$
(\lean{dtriPiece}). Every $c_i$ has two neighbours on the triangle, hence at most one parent outside it. Two distinct
triangle vertices cannot share an outside parent $w$: with the arc between them, $w$ would complete a transitive
triangle. So outside parents are private, and \cref{lem:Bdir} applies.
\end{outline}

When the underlying graph is triangle-free, Theorem~U3 follows from Theorem~N3 and \cref{lem:fewer-arcs}; its new
content is directed triangles. The fourth paper's Theorem~S~\cite[Theorem~6.13]{FourthPaper2026} proves the half bound
for scalar linear codes on a larger class: its condition ``no recovery triangle'' is TT-freeness, and its reciprocity
condition asks that a vertex with three parents have no child that is not a parent, which holds when the underlying
graph is subcubic. Theorem~U3 is thus the all-code version of Theorem~S on relations with subcubic underlying graph;
the all-code version of Theorem~S on its whole class is open.

\paragraph{Why it matters.} Directed triangles do not help a code. This matters for~(T) because native cone relations
can contain linearly many disjoint directed $3$-cycles~\cite[Proposition~3.10]{DirectedRecovery2026}, which rule out the
symmetric hypothesis of Theorem~N3 but not the hypotheses of this section.

\subsection{Acyclic high in-degree: Theorem T}\label{sec:T}

\begin{theorem}[Theorem T]\label{thm:T}
\formal{TheoremT.lean}{DCyc, exists_hasDCyc_of_sourceless, DCyc.chordless, DCyc.acyclic_of_missing, singlePiece, i5Piece, block_decomp, rateBound_thmT, card_sq_le_thmT}
Let $N$ be loopless and TT-free, and suppose that the set $\{v:\card{N(v)}\ge3\}$ of vertices with at least three
parents is acyclic for $N$. Then $N$ satisfies \eqref{eq:star}, and every storage code $C$ for $N$ over a nonempty
finite alphabet satisfies $\card C^2\le\card\alpha^{\card V}$. No bound on in- or out-degrees is assumed.
\end{theorem}

\begin{outline}
By strong induction on the set $R$ of remaining vertices (\lean{block_decomp}, with one block), in-degrees being counted
inside $R$:
\begin{enumerate}[label=(\alph*)]
\item If some $v\in R$ has at most one parent in $R$, then $\{v\}$ is a piece, its remaining parent being its outside
parent (\lean{singlePiece}, \cref{lem:Bdir} with $A=\{v\}$ and $B=\emptyset$).
\item Otherwise let $W_R$ be the set of vertices with exactly two parents in $R$. If $W_R$ is not acyclic, it contains a
nonempty set in which every vertex has a parent, hence a directed cycle (\lean{exists_hasDCyc_of_sourceless}). A
shortest directed cycle $C$ in $W_R$ (\lean{DCyc}) is induced (\lean{DCyc.chordless}): the only parent of each
$c\in C$ on $C$ is its predecessor. Hence each $c\in C$ has exactly one further parent in $R$, off $C$, its outside
parent $p(c)$. Fix $c_0\in C$ and let $A=\{c\in C:p(c)=p(c_0)\}$ and $B=C\setminus A$. If $c_0\to c_1$ on $C$, then
$p(c_1)\ne p(c_0)$, since otherwise $p(c_0)\to c_0\to c_1$ with $p(c_0)\to c_1$ is a transitive triangle. So $A$ and
$B$ are proper subsets of an induced cycle and hence acyclic (\lean{DCyc.acyclic_of_missing}), and every outside vertex
has its children on $C$ inside $A$ (if it is $p(c_0)$) or inside $B$. \Cref{lem:Bdir} applies (\lean{i5Piece}).
\item Otherwise $W_R$ is acyclic, and so is the set of vertices with at least three parents in $R$, a subset of the
acyclic set of the hypothesis. These two sets partition $R$, which is a piece by \cref{lem:Bdir}.
\end{enumerate}
Theorem~D concludes.
\end{outline}

Strongly connected components can be treated separately, since a parent in an earlier component can be contracted. The
formal statement is a block form.

\begin{corollary}[Corollary $\Tint$]\label{cor:Tint}
\partformal{TheoremT.lean}{TTFreeOn, BlockOrder, blocks_decomp, rateBound_blocks}{the existence of a topological order of
the strongly connected components, which is standard; it is proved below.}
Let $N$ be loopless. Suppose that no strongly connected component of $N$ contains a transitive triangle, and that the set
$\Tint$ of vertices with at least three parents inside their own strongly connected component is acyclic. Then $N$
satisfies \eqref{eq:star}. There is no restriction on parents outside a vertex's component.
\end{corollary}

The formal theorem (\lean{rateBound_blocks}) is the following block form: if $N$ is loopless and $V$ is the union of
pairwise disjoint blocks
$S_1,\dots,S_r$ such that every parent of a vertex of $S_j$ lies in $S_1\cup\dots\cup S_j$, no block contains a
transitive triangle, and each set $\{v\in S_j:\card{N(v)\cap S_j}\ge3\}$ is acyclic, then $N$ satisfies
\eqref{eq:star}.

\begin{outline}
Induction over the list of blocks (\lean{blocks_decomp}): each block $S_j$ is decomposed by the procedure of Theorem~T
run inside $S_j$ (\lean{block_decomp}), with the earlier blocks in the context. Parents in earlier blocks are then free,
so in-degrees in~(a)--(c) are counted inside the remaining part of $S_j$, and TT-freeness is needed only inside $S_j$
(\lean{TTFreeOn}).
\end{outline}

\begin{informalproof}
We derive the corollary from the block form. The \emph{condensation} of $N$ has the strongly connected components as vertices and an arc $S\to S'$, $S\ne S'$,
whenever some arc of $N$ goes from $S$ to $S'$. It has no directed cycle, since components on such a cycle would be
strongly connected to each other. Hence it has a topological order: repeatedly remove a component that receives no arc
from the remaining ones. Listing the components $S_1,\dots,S_r$ in this order, every parent of a vertex of $S_j$ lies in
$S_1\cup\dots\cup S_j$. The set $\{v\in S_j:\card{N(v)\cap S_j}\ge3\}$ is $\Tint\cap S_j$, which is acyclic because
acyclicity passes to subsets. The block form applies.
\end{informalproof}

When $\Tint=\emptyset$, that is, when every vertex has at most two parents inside its component, Corollary~$\Tint$ is
the TT-free case of the two-input theorem of \emph{Directed Recovery
Structures}~\cite[Theorem~5.1]{DirectedRecovery2026}: TT-free components contain no bidirected triangles, so that
theorem's correction term vanishes. The corollary strictly contains this case, because it allows vertices with three or
more internal parents as long as they induce an acyclic sub-relation. Its proof dismantles a bidirected odd cycle by
contracting one digon, a shortest directed cycle of in-degree-two vertices, and does not use the directed Brooks
theorem on which the proof in~\cite{DirectedRecovery2026} relies. Conversely, \cite[Theorem~5.1]{DirectedRecovery2026}
also covers components that contain transitive triangles, which we do not.

\paragraph{Why it matters.} Theorem~T and Corollary~$\Tint$ have no degree bound. For~(T) they locate where a coding
gain on TT-free side information must come from: from directed cycles through vertices that have at least three parents
in their own component. In a cone relation, after the middle vertices of transitive triangles have been deleted at the
price of \cref{cor:paid}, a construction needs such cycles on a large fraction of the coordinates
(\cref{sec:construction}).

\subsection{Blocks with private outside parents: Theorem PC}\label{sec:PC}

The pieces of Theorems~N3, U3 and~T are small. The next theorem takes whole blocks as pieces, in any order; a block
then needs only an internal certificate.

\begin{definition}[Elimination certificates of a block]\label{def:elimon}
\formal{Blocks.lean}{ElimOn}
Let $K\subseteq V$. An \emph{elimination certificate} of $N[K]$ consists of an elimination list of $K$ without
repetitions, records partitioning $K$, and own and supplying records $a(v)\ne b(v)$ with $v\in P_{a(v)}$, such that at
every step $t$ the parents of $v_t$ inside $K$ lie in $S_t\cup P_{b(v_t)}$.
\end{definition}

This is a global elimination certificate of the induced relation $N[K]$ (\cref{def:elimcert}), presented by its list.

\begin{lemma}[Internal certificates give pieces]\label{lem:elimon-piece}
\formal{Blocks.lean}{ElimOn.piece}
Let $N[K]$ have an elimination certificate. Suppose that every vertex of $K$ has at most one parent outside $K$, and that
every vertex outside $K$ is a parent of at most one vertex of $K$. Then $K$ has a piece certificate in every context $U$
disjoint from $K$.
\end{lemma}

\begin{outline}
Keep the list and the records. The outside parent of $v\in K$ is its parent outside $U\cup K$, if any; its other
parents lie in $U$ or in $K$, where the internal certificate applies. Outside parents are private, so
\cref{lem:cosupply} gives nesting. (The formal statement asks for at most one parent in $V\setminus(U\cup K)$ only.)
\end{outline}

\begin{lemma}[Positional certificates; chordless digon cycles]\label{lem:cycle-cert}
\formal{Circulant.lean}{elimOn_ofFn, elimOn_cycle}
Let $K=\{k_0,\dots,k_{m-1}\}$ be nonempty and listed without repetitions.
\begin{enumerate}[label=(\alph*)]
\item\label{it:positional} Let $\mathrm{col}$ and $\mathrm{bcol}$ be colourings of the positions $0,\dots,m-1$ with
$\mathrm{col}(j)\ne\mathrm{bcol}(j)$ for every $j$. If $k_i\to k_j$ implies $i<j$ or $\mathrm{col}(i)=\mathrm{bcol}(j)$,
then $N[K]$ has an elimination certificate (the formal proof uses the list $k_0,\dots,k_{m-1}$ and the colour classes as
records).
\item\label{it:cycle} If $m\ge4$ and every arc inside $K$ joins cyclically consecutive vertices ($k_i\to k_j$ only if
$i\equiv j\pm1\pmod m$), then $N[K]$ has an elimination certificate.
\end{enumerate}
\end{lemma}

\begin{outline}
\ref{it:positional} The own record of $k_j$ is its colour and its supplying record is $\mathrm{bcol}(j)$; a parent $k_i$
that comes later has colour $\mathrm{bcol}(j)$. \ref{it:cycle} Use the colouring of the cycle pieces of Theorem~N3, with
$\mathrm{bcol}(j)=\mathrm{col}(j+1)$ for $j<m-1$ and $\mathrm{bcol}(m-1)=0$.
\end{outline}

Along a list $k_0,\dots,k_{L-1}$, call an arc $k_u\to k_v$ with $u\ge v+3$ a \emph{downward chord}.

\begin{lemma}[Lemma S: chorded paths]\label{lem:S}
\formal{PathCert.lean}{DownChord, pathRoot, pathRoot_succ_ne, elimOn_path}
Let $K=\{k_0,\dots,k_{L-1}\}$ be nonempty and listed without repetitions. Suppose that every arc $k_i\to k_j$ inside
$K$ joins consecutive positions, $\card{i-j}=1$, or skips at least two positions, $\card{i-j}\ge3$, and that every vertex
of $K$ is the head of at most one downward chord and the tail of at most one. Then $N[K]$ has an elimination certificate
(the formal proof uses the list $k_0,\dots,k_{L-1}$ and, as records, the classes of the equivalence relation on positions
generated by the \emph{links} $u\sim v+1$, one for each downward chord $k_u\to k_v$, together with one empty
record).
\end{lemma}

\begin{outline}
When $k_j$ is eliminated, its parents inside $K$ that come later are $k_{j+1}$ (if $j<L-1$) and the tail $k_u$ of the
downward chord to $k_j$, if there is one; these two are linked, so they lie in one class, which serves as supplying record
(through \cref{lem:cycle-cert}\ref{it:positional}). It remains to show that $j$ and $j+1$ are never in the same class
(\lean{pathRoot_succ_ne}). A link joins $u$ to $v+1\le u-2$. A position is the upper end of at most one link, from the
downward chord out of it, and the lower end of at most one link, from the downward chord into the position below it,
because these chords are unique. So the classes are chains whose positions decrease by at least two at each link:
following links downwards from any position reaches a lowest position, its \emph{root} (\lean{pathRoot}), and two
positions of one class differ by at least two.
\end{outline}

A \emph{chorded digon path} is a set $K=\{k_0,\dots,k_{L-1}\}$ such that $k_i\leftrightarrow k_{i+1}$ is a digon for
$0\le i<L-1$ and the other arcs inside $K$, the \emph{chords}, number at most one into and at most one out of every vertex.
If $N$ is loopless and $N[K]$ is TT-free, a chorded digon path satisfies the hypotheses of Lemma~S
(\lean{ChordedPath.elimOn}): a chord between $k_i$ and $k_{i\pm2}$ would close a transitive triangle with the two digons
around $k_{i\pm1}$.

\begin{proposition}[Proposition CC: chorded circulant cycles]\label{prop:CC}
\formal{Circulant.lean}{circCol, circB, elimOn_circulant, circulant_sweep}
Let $3\le k\le m-3$ and let $K=\{k_0,\dots,k_{m-1}\}$ be listed without repetitions so that every arc inside $K$ has the
form $k_{j\pm1}\to k_j$ or $k_{j-k}\to k_j$, indices modulo $m$. Put $\delta=m-1-k$.
\begin{enumerate}[label=(\alph*)]
\item\label{it:sweep} If $\delta$ divides neither $m-1$ nor $m-3$, then $N[K]$ has an elimination certificate (the
formal proof uses the sweep along $k_0,k_1,\dots,k_{m-1}$ described below).
\item\label{it:arith} If $2k\not\equiv\pm1\pmod m$, then the condition of~\ref{it:sweep} holds for the offset $k$ or for
the offset $m-k$.
\end{enumerate}
\end{proposition}

\begin{outline}
\ref{it:sweep} (\lean{elimOn_circulant}, through \cref{lem:cycle-cert}\ref{it:positional}.) Position $0$ has a record of
its own; a position $p\ne0$ gets the record of its residue modulo $\delta$, except that the residue class of $m-1$ is merged
into that of $1$ (\lean{circCol}). The supplying record of $p<m-1$ is the record of $p+1$, and that of $m-1$ is the record
of position~$0$ (\lean{circB}). The later parents of $k_p$ are $k_{p+1}$; for $p=0$ also $k_{m-1}$, which has the merged
colour of $k_1$; and, for $p<k$, the tail $k_{p-k+m}$ of the chord, whose position $p-k+m=(p+1)+\delta$ is congruent to
$p+1$. Own and supplying records differ: for $1\le p<m-1$ the residues of $p$ and $p+1$ differ, and the merge identifies
them only if $p\equiv m-1$ and $p+1\equiv1$, which forces $\delta\mid m-1$, or $p\equiv1$ and $p+1\equiv m-1$, which
forces $\delta\mid m-3$.
\ref{it:arith} (\lean{circulant_sweep}.) For the offset $m-k$ the role of $\delta$ is played by $\delta'=k-1$, and
$m-1=\delta+\delta'+1$, $m-3=\delta+\delta'-1$. If the condition fails for both offsets, then $\delta\mid\delta'\pm1$ and
$\delta'\mid\delta\pm1$ with $\delta,\delta'\ge2$, which forces $\card{\delta-\delta'}=1$ and then $2k=m\pm1$.
\end{outline}

\emph{With the mirror sweep, if $2k\not\equiv\pm1\pmod m$, then $N[K]$ has an elimination certificate.} Indeed, by
part~\ref{it:arith} the condition of part~\ref{it:sweep} holds for the offset $k$, and part~\ref{it:sweep} applies, or
for the offset $m-k$. In the second case relabel $k'_j=k_{-j}$: the relabelling $j\mapsto-j$ turns the arcs
$k_{j\pm1}\to k_j$ into arcs of the same form and the chords $k_{j-k}\to k_j$ into chords $k'_{j-(m-k)}\to k'_j$ of
offset $m-k$, so part~\ref{it:sweep} applies to the list $k'_0,k'_1,\dots,k'_{m-1}$, that is, to $k_0,k_{m-1},\dots,k_1$.
This is the \emph{mirror sweep}; it is not a single Lean declaration, and the formal proof of Corollary~CC below realizes
it by indexing cosets with powers of $s^{-1}$. The condition $2k\not\equiv0$, which makes the chords one-way, is not
needed. In a TT-free circulant $\Cay(\Z_m,\{1,-1,k\})$ with one-way chords the hypotheses
hold, since chords of length $2$ and chords with $2k\equiv\pm1$ close transitive triangles.

\begin{theorem}[Theorem PC: blocks with private outside parents]\label{thm:PC}
\formal{PathCert.lean}{ChordedPath, ChordedPath.elimOn, rateBound_pc}
\formal{Blocks.lean}{ChordlessSubcubic, localGraph, chordless_block_decomp}
\formal{N3.lean}{block_decomp_of_cycles}
Let $N$ be loopless and let $V$ be partitioned into blocks. Suppose:
\begin{enumerate}[label=(P\arabic*),leftmargin=3em]
\item\label{it:P1} every vertex has at most one parent outside its block;
\item\label{it:P2} every vertex is a parent of at most one vertex of each other block;
\item\label{it:P3} every block $K$ is
  \begin{enumerate}[label=(\alph*)]
  \item \emph{chordless subcubic}: every arc inside $K$ belongs to a digon, every vertex of $K$ has at most three
  parents, and $N[K]$ is TT-free; or
  \item a chorded digon path with $N[K]$ TT-free; or
  \item such that $N[K]$ has an elimination certificate.
  \end{enumerate}
\end{enumerate}
Then $N$ satisfies \eqref{eq:star}; consequently every storage code $C$ for $N$ over a nonempty finite alphabet satisfies
$\card C^2\le\card\alpha^{\card V}$ (\cref{thm:codes}). Out-degrees are not restricted.
\end{theorem}

\begin{outline}
By Lemma~S, blocks of kind~(b) have elimination certificates and are therefore of kind~(c)
(\lean{ChordedPath.elimOn}). The blocks are taken as pieces in an arbitrary order (\lean{rateBound_pc}), each with the earlier blocks as context.
By~\ref{it:P1} and~\ref{it:P2}, every vertex of a block $K$ has at most one parent outside $K$, and outside parents are
private. A block of kind~(c) is a piece by \cref{lem:elimon-piece}. A block of kind~(a) is decomposed relative to its
context by the shortest-cycle procedure of Theorem~N3 (\lean{chordless_block_decomp}, \lean{block_decomp_of_cycles}),
applied to its local graph (\lean{localGraph}), which joins each vertex of $K$ to its parents. Inside $K$ the children of
a vertex are its digon partners, so every vertex of $K$ has local degree at most three; a vertex outside $K$ has at most
one local neighbour, by~\ref{it:P2}; and the local graph has no triangle inside $K$, since such a triangle would consist
of digons and contain a transitive triangle. The cycle pieces and the final relative $2$-colouring then have private
outside parents.
\end{outline}

The blocks are usually far from minimum pieces; Theorem~PC is not a shortest-cycle argument. Its natural blocks are the
components of the digon graph.

\begin{corollary}[Theorem PC for digon components]\label{cor:PC-digon}
\notformalized
Let $N$ be loopless and TT-free. Suppose that
\begin{enumerate}[label=(H\arabic*),leftmargin=3em]
\item every vertex has at most one one-way parent, and
\item no vertex has one-way arcs to two vertices of the same component of the digon graph $G_2(N)$,
\end{enumerate}
and that every component $K$ of $G_2(N)$ has no one-way arc inside it and only vertices with at most three parents, or
induces a path of $G_2(N)$, or is such that $N[K]$ has an elimination certificate (for instance a cycle of $G_2(N)$ of length $m$ whose
one-way arcs inside $K$ all have the same offset $k$ with $2k\not\equiv\pm1\pmod m$, by \cref{prop:CC} with the mirror
sweep; TT-freeness gives $3\le k\le m-3$). Then $N$ satisfies \eqref{eq:star}.
\end{corollary}

\begin{informalproof}
Take the components of $G_2(N)$ as blocks. The digon partners of a vertex lie in its own component, so its parents
outside its component are one-way; by~(H1) there is at most one, which is~\ref{it:P1}. A vertex outside a component $K$
has only one-way arcs into $K$, at most one by~(H2), which is~\ref{it:P2}. For~\ref{it:P3}: a component without one-way
arcs inside it has every internal arc in a digon, and with the in-degree bound and TT-freeness it is of kind~(a). If $K$
induces a path $k_0,\dots,k_{L-1}$ of $G_2(N)$, its consecutive pairs are digons, and every other arc inside $K$ is
one-way, because a digon inside $K$ is an edge of $G_2(N)[K]$, which is the path. Each vertex is the head of at most one
one-way arc, by~(H1), and the tail of at most one one-way arc into $K$, by~(H2); so $K$ is a chorded digon path, of
kind~(b). The remaining components are of kind~(c) by hypothesis. Theorem~PC applies.
\end{informalproof}

Every chordless cycle of $G_2(N)$ is of kind~(a), and also of kind~(c) by \cref{lem:cycle-cert}\ref{it:cycle}; it has
length at least four, since $G_2(N)$ is triangle-free. Whether every TT-free chorded digon cycle is of kind~(c) is
\cref{conj:K}.

\paragraph{Why it matters.} Theorem~PC covers relations that none of Theorems~N3, U3 and~T covers, and puts no bound on
out-degrees: a vertex may send one one-way arc into each of arbitrarily many blocks. For~(T), side information built
from two-way exchanges inside blocks, linked by one-way arcs that are private in the sense of~\ref{it:P1}
and~\ref{it:P2}, gives no coding gain. With the digon components as blocks, a construction needs, on a large scale,
vertices that receive several one-way arcs or send several into one component, or components without certificates. The
relation $\Ex$ of \cref{rem:X15} has vertices of the first kind, and yet it has a global certificate.

\subsection{Cayley relations: Corollary CC}\label{sec:CC}

\begin{corollary}[Corollary CC: Cayley relations]\label{cor:CC}
\formal{Cayley.lean}{cayN, coset_steps, mem_coset_iff, rateBound_cayley, card_sq_le_cayley}
Let $G$ be a finite group and let $s,t\in G$ be such that $S=\{s,s^{-1},t\}$ has three elements and does not contain
$1$. Let $N$ be the Cayley relation with arcs $g\to gx$ for $x\in S$, so $N(q)=\{qs^{-1},qs,qt^{-1}\}$. If $N$ is
TT-free, then $N$ satisfies \eqref{eq:star}, and every storage code $C$ for $N$ over a nonempty finite alphabet
satisfies $\card C^2\le\card\alpha^{\card G}$.
\end{corollary}

\begin{outline}
The relation is loopless because $1\notin S$. If $t^2=1$, every arc is reciprocated, the underlying graph has degree at
most three, and Theorem~U3 applies. Otherwise take the left cosets $g\langle s\rangle$ as blocks. The order $m$ of $s$ is at least four: $m\ne1,2$ because $s\ne1$ and $s\ne s^{-1}$, and $m\ne3$
because $g\to gs\to gs^2=gs^{-1}$ with $g\to gs^{-1}$ would be a transitive triangle. The only parent of $q$ outside its
coset is $qt^{-1}$, and the only possible child outside it is $qt$; so~\ref{it:P1} and~\ref{it:P2} hold. Index the coset of
$g$ by $j\mapsto gs^j$ (\lean{mem_coset_iff}); in-neighbours through $s^{\pm1}$ then sit at neighbouring positions
(\lean{coset_steps}). If $t\notin\langle s\rangle$, its internal arcs join consecutive positions,
and \cref{lem:cycle-cert}\ref{it:cycle} applies. If $t=s^k$, its internal arcs are $gs^{j\pm1}\to gs^j$ and
$gs^{j-k}\to gs^j$. Here $k\ne0,\pm1$ because $\card S=3$ and $1\notin S$; $k\ne\pm2$ and $2k\not\equiv\pm1$ because
each would close a transitive triangle. So $3\le k\le m-3$ and \cref{prop:CC} applies, with the generator $s$ or, for the
mirror sweep, $s^{-1}$ and the offset $m-k$. Theorem~PC concludes.
\end{outline}

If instead $S$ has at most two elements and does not contain~$1$, every vertex has at most two parents, and Theorem~T
applies when $N$ is TT-free. For example,
$\Cay(\Z_9,\{1,-1,3\})$, with $N(q)=\{q-1,q+1,q-3\}$, is covered by Corollary~CC (one block, $m=9$,
$k=3$, $\delta=5$). It is covered by none of the earlier results: its underlying graph is $4$-regular, so neither
Theorem~N3, even with \cref{lem:fewer-arcs}, nor Theorem~U3 applies; it is strongly connected with every in-degree equal
to three, so Theorem~T and Corollary~$\Tint$ do not apply; and it is not $2$-dicolourable, since the two ends of a
digon must lie in different acyclic sets and its digons form a $9$-cycle. It also shows that the private pieces of Theorem~N3 can fail to exist (\cref{prop:no-private}).

\paragraph{Why it matters.} Cayley relations are the side information of uniform, translation-invariant schedules. For
every finite group, the pattern ``two-way exchange along one generator, one-way transfer along another'' gives no
coding gain at any size. For~(T), symmetric constructions of this shape are excluded, whatever the group.
